\section{Methodology}
\label{sec:method}

We designed a generative evaluation pipeline to extract and analyze the linguistic features of LLM-generated spatial narratives. Unlike traditional benchmarks that ask a model to "solve" a pathfinding problem, our approach provides the model with the ground-truth layout and the correct path, then asks it to "describe" the journey.

\para{Dataset and Environment}
We sampled 50 layouts and traversal questions from the \babitask dataset \cite{weston2015towards}. Each layout defines a set of rooms and their topological relationships (e.g., "The kitchen is north of the hallway"). We extracted the optimal path between the start and end points specified in each task.

\para{Models and Prompting Conditions}
We evaluated two state-of-the-art models from OpenAI: \gpto and \gptom. For each traversal task, we generated narratives under two conditions:
\begin{itemize}[leftmargin=*,itemsep=0pt,topsep=0pt]
    \item \standardprompt: "Write a short narrative (3-5 sentences) describing your journey from [start] to [end]."
    \item \humanprompt: "Write a creative, highly varied, and human-like story (3-5 sentences)... Avoid robotic, formulaic sentence structures."
\end{itemize}
These conditions allow us to distinguish between the model's default "learned" spatial grammar and its ability to simulate naturalistic human descriptions when explicitly instructed.

\para{Linguistic Feature Extraction}
We processed the generated narratives using \textsc{spaCy} and \textsc{NLTK} to extract several key linguistic features:
\begin{itemize}[leftmargin=*,itemsep=0pt,topsep=0pt]
    \item {\bf \prepentropy}: We calculated the Shannon entropy of the distribution of spatial prepositions (Adpositions) used in each narrative. Higher entropy indicates a more varied and descriptive spatial vocabulary.
    \item {\bf \sentvar}: We measured the variance in sentence length (word count) within each narrative. Low variance indicates a repetitive, "robotic" rhythmic structure.
    \item {\bf POS Bigram Distributions}: We computed the frequency of Part-of-Speech bigrams to identify recurring syntactic structures.
\end{itemize}

\para{Baseline Comparison}
Due to the lack of a large-scale, controlled dataset of human-written spatial narratives for these specific layouts, we use the \humanprompt condition as a proxy for the model's "naturalistic" upper bound and compare the \standardprompt outputs against it. We also compare the distributions between \gpto and \gptom to evaluate the impact of model scale.
